%! Polynomial Functions and End Behavior — Unit 2, Lesson 2.5 — Group Activity
\documentclass[10pt]{article}
\usepackage{precalculus-article}
\usepackage{precalculus-boxes}
\newcommand{\TallMath}[1]{$\displaystyle #1\rule[-1.4em]{0pt}{3.2em}$}

\begin{document}

\pageheader{Unit 2, Lesson 2.5}{Group Activity}

\vspace{0.12in}

\begin{scenariobox}[The End Behavior Desk]{plumlight}
\small
Your group takes the end-behavior desk: every polynomial that arrives leaves
with both of its far ends described, in limit notation, without anyone graphing
anything. Tier E reads the two pre-drawn graphs below --- no tier draws
anything.

\begin{center}
\begin{tikzpicture}[x=0.7cm, y=0.28cm]
  \draw[linegray, very thin, xstep=1, ystep=2] (-2.4,-5) grid (2.4,5);
  \draw[->, thick] (-2.8,0) -- (2.9,0) node[below right] {\small $x$};
  \draw[->, thick] (0,-5.6) -- (0,5.6) node[above] {\small $y$};
  \draw[very thick, plum] plot [smooth] coordinates
    {(-2.2,-4.07) (-2,0) (-1.7,3.21) (-1.4,4.0) (-1,3) (-0.5,0.94)
     (0,0) (0.5,0.94) (1,3) (1.4,4.0) (1.7,3.21) (2,0) (2.2,-4.07)};
  \node[below] at (0,-6.2) {\small \textbf{Graph P}};
\end{tikzpicture}
\hspace{1.6cm}
\begin{tikzpicture}[x=0.7cm, y=0.28cm]
  \draw[linegray, very thin, xstep=1, ystep=2] (-2.4,-5) grid (2.4,5);
  \draw[->, thick] (-2.8,0) -- (2.9,0) node[below right] {\small $x$};
  \draw[->, thick] (0,-5.6) -- (0,5.6) node[above] {\small $y$};
  \draw[very thick, plum] plot [smooth] coordinates
    {(-2.2,4.05) (-2,2) (-1.8,0.43) (-1.5,-1.13) (-1,-2) (-0.5,-1.38)
     (0,0) (0.5,1.38) (1,2) (1.5,1.13) (1.8,-0.43) (2,-2) (2.2,-4.05)};
  \node[below] at (0,-6.2) {\small \textbf{Graph Q}};
\end{tikzpicture}
\end{center}
\end{scenariobox}

\vspace{0.1in}

\boxguard
\begin{tcolorbox}[colback=white, colframe=black!40,
    title=\textbf{Tier R --- Remediate},
    fonttitle=\bfseries, arc=2mm, left=3mm, right=3mm, top=2mm, bottom=2mm]

\begin{enumerate}[itemsep=8pt]
  \item The leading term is already done for you. Fill the rest of each row
        from the four-case table in your notes.

        \smallskip
        \renewcommand{\arraystretch}{1.6}
        \begin{center}
        \small
        \begin{tabular}{|c|c|c|c|c|}
        \hline
        \textbf{Leading term} & \textbf{Even or odd?} & \textbf{Sign of $a_n$?} & $\lim_{x \to \infty} p(x)$ & $\lim_{x \to -\infty} p(x)$ \\
        \hline
        $4x^2$   & \blank{1.6cm} & \blank{2cm} & \blank{1.4cm} & \blank{1.4cm} \\
        \hline
        $-x^5$   & \blank{1.6cm} & \blank{2cm} & \blank{1.4cm} & \blank{1.4cm} \\
        \hline
        $-3x^8$  & \blank{1.6cm} & \blank{2cm} & \blank{1.4cm} & \blank{1.4cm} \\
        \hline
        $7x^3$   & \blank{1.6cm} & \blank{2cm} & \blank{1.4cm} & \blank{1.4cm} \\
        \hline
        \end{tabular}
        \end{center}

  \item Complete the summary sentence.

        \begin{enumerate}[label=(\alph*), itemsep=4pt]
          \item An \textbf{even} degree sends both ends in the \blank{2.2cm}
                direction.
          \item An \textbf{odd} degree sends the ends in \blank{2.4cm}
                directions.
          \item A \textbf{negative} leading coefficient \blank{1.8cm} both
                ends.
        \end{enumerate}

  \item Two of the four polynomials in question 1 have a graph that must cross
        the $x$-axis. Which two? \quad \blank{4.4cm}
\end{enumerate}
\end{tcolorbox}

\vspace{0.1in}

\boxguard
\begin{tcolorbox}[colback=white, colframe=black!40,
    title=\textbf{Tier A --- Approaching Proficiency},
    fonttitle=\bfseries, arc=2mm, left=3mm, right=3mm, top=2mm, bottom=2mm]

\begin{enumerate}[itemsep=8pt]
  \item Find the leading term yourself, then write both limit statements.
        Read the exponents, not the order the terms appear in.

        \begin{enumerate}[label=(\alph*), itemsep=6pt]
          \item $f(x) = 3x^4 - 9x^2 + 1$

                Leading term: \blank{2cm} \quad
                $\lim_{x \to \infty} f(x) = $ \blank{1.6cm} \quad
                $\lim_{x \to -\infty} f(x) = $ \blank{1.6cm}

          \item $g(x) = 12 - 5x^3 + 2x^6 - x^7$

                Leading term: \blank{2cm} \quad
                $\lim_{x \to \infty} g(x) = $ \blank{1.6cm} \quad
                $\lim_{x \to -\infty} g(x) = $ \blank{1.6cm}
        \end{enumerate}

  \item Now in factored form. Add the exponents for the degree, multiply the
        leading pieces for the leading term --- do not expand.

        \begin{enumerate}[label=(\alph*), itemsep=6pt]
          \item $h(x) = (x-4)(x+1)(x-7)$

                \begin{work}
                  \text{degree of } h &= 1 + 1 + 1 \\
                                      &= 3 \\
                  \text{leading term} &= x \cdot x \cdot x \\
                                      &= x^{3}
                \end{work}

                $\lim_{x \to \infty} h(x) = $ \blank{1.6cm} \quad
                $\lim_{x \to -\infty} h(x) = $ \blank{1.6cm}

          \item $k(x) = -2(x+3)^2(x^2+1)$

                \begin{work}
                  \text{degree of } k &= 2 + 2 \\
                                      &= 4 \\
                  \text{leading term} &= -2 \cdot x^{2} \cdot x^{2} \\
                                      &= -2x^{4}
                \end{work}

                $\lim_{x \to \infty} k(x) = $ \blank{1.6cm} \quad
                $\lim_{x \to -\infty} k(x) = $ \blank{1.6cm}
        \end{enumerate}

  \item Which of the four functions above are guaranteed to cross the
        $x$-axis, and what is it about them that guarantees it?

        \vspace{0.05in}
        \writelines{2}
\end{enumerate}
\end{tcolorbox}

\vspace{0.1in}

\boxguard
\begin{tcolorbox}[colback=white, colframe=black!40,
    title=\textbf{Tier E --- Extension},
    fonttitle=\bfseries, arc=2mm, left=3mm, right=3mm, top=2mm, bottom=2mm]

\begin{enumerate}[itemsep=8pt]
  \item Read Graphs P and Q at the top of the page. Neither equation is given
        --- the ends are enough.

        \begin{enumerate}[label=(\alph*), itemsep=4pt]
          \item Graph P: \; $\lim_{x \to \infty} y = $ \blank{1.4cm},
                \; $\lim_{x \to -\infty} y = $ \blank{1.4cm}, \;
                so the degree is \blank{1.4cm} and $a_n$ is \blank{1.8cm}.
          \item Graph Q: \; $\lim_{x \to \infty} y = $ \blank{1.4cm},
                \; $\lim_{x \to -\infty} y = $ \blank{1.4cm}, \;
                so the degree is \blank{1.4cm} and $a_n$ is \blank{1.8cm}.
        \end{enumerate}

  \item For $p(x) = x^4 - 500x^3$, find the exact input past which the leading
        term is larger in size than the rest.

        \begin{work}
          \frac{-500x^3}{x^4} &= \frac{-500}{x} \\
          \left|\frac{-500}{x}\right| &< 1 \\
          x &> 500
        \end{work}

        The leading term takes over for $x > $ \blank{1.6cm}. Fill in the
        ratio at three inputs:

        \smallskip
        \renewcommand{\arraystretch}{1.5}
        \begin{center}
        \begin{tabular}{|c|c|c|c|}
        \hline
        $x$ & $50$ & $500$ & $5{,}000$ \\
        \hline
        rest $\div$ leading & \blank{1.6cm} & \blank{1.6cm} & \blank{1.6cm} \\
        \hline
        \end{tabular}
        \end{center}

  \item Write a polynomial in \textbf{factored form} that has even degree, a
        negative leading coefficient, and exactly three distinct real zeros.

        \vspace{0.05in}
        $p(x) = $ \blank{8cm}

  \item Explain why $p(x) = x^5 - 900x^4 + 12$ and $q(x) = x^5 + 7x$ have
        exactly the same end behavior, even though their graphs look nothing
        alike on a standard window.

        \vspace{0.05in}
        \writelines{3}
\end{enumerate}
\end{tcolorbox}

\end{document}
